\chapter{Coverings, equivariant AR theory and graded surfaces}
\label{cover:chapter}
The new covering calculation in Section~\ref{cover:section} is useful for
computations, precisely because its limitations are explicit. The parent
algebra is our infinite elliptic theta algebra, its truncation is a finite
quotient, and its cover has vertices indexed by integers. They are three
different categories. The nine-dimensional module in
Proposition~\ref{cover:band} and the higher-genus family in
Corollary~\ref{cover:higher} show that even this genuine covering does not
reach every module or bounded derived object.

\section{Reading the three new sources together}
The lecture \cite{CoveringLecture2026} supplies the covering language:
locally bounded categories, exact push-down/pull-up and graded adjunctions.
Its Theorems~4.23 and~4.28 give precoverings on the specified module and bounded
derived categories. The AR-quiver covering theorem~4.18 assumes an actual
Galois covering between Hom-finite Krull--Schmidt categories; density cannot
be dropped. A universal bound-quiver cover also depends on its relations and
presentation. It is not automatically a covering of the compact surface.

The GLS paper \cite{TameGLS2026} works over an algebraically closed
characteristic-zero field. Its finite-dimensional GLS algebras are not
automatically our surface orders. Equivariantization transports string/band
information to certain non-string algebras. Theorem~B includes
$\Z D_\infty$, $\Z A_\infty^\infty$ and tube components, with its stated
indexing. Theorem~C corrects a naive root-only prediction by the additional
rank vectors $\beta_i+(2r+1)\delta$ in the named tame affine family.
Here vertexwise local freeness and its rank vector concern the GLS vertex
algebras; they are not freeness over the central DVR of a surface order.

Schroll's survey \cite{SchrollGeometric2026}, particularly Sections~4--5,
organizes graded gentle categories through curves, local systems, crossings
and twisted complexes, with grading encoded by a line field. Its framework
distinguishes perfect objects from finite-total-cohomology objects, and
properness from homological smoothness. These distinctions matter for our
infinite algebras and for any completed or operator-algebra version.

\section{A cover that really preserves the relations}
At finite $L$ the cover is locally bounded: only lengths below $L$ survive.
Without the cutoff, each vertex has infinitely many nonzero outgoing Hom
spaces, so the locally bounded hypothesis used above fails. Locally finite
dimensional Hom spaces alone do not repair that failure. Conversely the
truncated algebra has extra length relations and is not entitled to all
quadratic-gentle surface-model theorems.

\begin{center}
\begin{tikzpicture}[>=Stealth,scale=.92,every node/.style={font=\small}]
\foreach \r in {0,1,2,3,4,5}{
  \node (v\r) at (1.3*\r,1.8) {$\r$};
}
\foreach \r/\s in {0/1,1/2,2/3,3/4,4/5}{
  \draw[->,bend left=22] (v\r) to node[above] {$a$} (v\s);
  \draw[->,bend right=22] (v\r) to node[below] {$b$} (v\s);
}
\node (e0) at (0,0) {$e_0$};
\node (e1) at (2.6,0) {$e_1$};
\node (e2) at (5.2,0) {$e_2$};
\draw[->,bend left=18] (e0) to node[above] {$a_0,b_0$} (e1);
\draw[->,bend left=18] (e1) to node[above] {$a_1,b_1$} (e2);
\draw[->,bend left=30] (e2) to node[below] {$a_2,b_2$} (e0);
\foreach \r/\s in {0/0,1/1,2/2,3/0,4/1,5/2}{
  \draw[densely dotted,->,gray] (v\r) to (e\s);
}
\node[anchor=west] at (0,2.7) {Cover: $r\mapsto r+3$ is a deck shift};
\end{tikzpicture}
\par\smallskip{\small Mixed paths and all paths of length $L$ vanish on both rows.}
\end{center}
The dotted lines are the residue map, not irreducible module maps. The top
row is a finite window of the quiver of the cover, not an AR component.
The proof of the Hom decomposition, the unequal Hom-control defects and the
full nine-dimensional calculation appear in
Sections~\ref{cover:section} and~\ref{cover:details}; they are not inferred
from the drawing. The genus-four instance has dimension~27, and the
genus-nine instance dimension~57, with the same exact lift obstruction.

\section{Equivariant modules, not an unproved folding rule}
The branch swap preserves $J^L$. Thus it acts on $B_L$ and defines the
skew-group algebra $B_L\# k[\Z/2]$. An equivariant object needs an
isomorphism $c:M\to\sigma M$ satisfying the cocycle equation
$\sigma(c)c=1$. A fixed isomorphism class alone does not choose $c$.
For the nine-dimensional witness, the unequal branch ranks rule out even
$M\simeq\sigma M$; the induced pair $M\oplus\sigma M$ instead has a
canonical swapping action. This is a concrete use of equivariant machinery
without asserting a GLS identification or drawing a folded AR component
whose almost-split property has not been checked.

\section{Derived and symplectic data needed for a determinant}
Let $\operatorname{Perf}(A)$ be the compact objects of $D(A)$ and write
$D_{\rm fd}(A)$ for finite-dimensional total cohomology. For our infinite
zero-graded $A$, its regular projective is perfect but has infinite
$k$-dimension; hence it does not belong to $D_{\rm fd}(A)$. This is already
a counterexample to silently embedding the former category in the latter.
In proper smooth finite-dimensional situations the familiar equivalences
have stronger conclusions; these hypotheses cannot be omitted here.

A curve model must retain the dissection, stops or punctures, line field,
intersection grading, local systems and dg or $A_\infty$ enhancement.
For perfect complexes, the finite deck decomposition of the Hom complex is
proved in Section~\ref{cover:details}. For a matrix of paths representing a
differential one must still check $d^2=0$; intersections without degrees and
composition data do not prove this identity. A marked surface used for the
Fukaya model is not automatically the original compact Riemann surface
with its analytic metric and all boundary data forgotten.

Completion to $\Lambda$ requires continuous modules and the chosen topology.
The finite quotient modules extend continuously because a power of $J$
annihilates them. They are not central-DVR-free lattices: in the
nine-dimensional example the central length-three parameter already acts
by zero. The cutoffs $A/J^L$ must not be renamed $\Lambda/u^L$.

For a useful arithmetic determinant, these categorical constructions still
need a prime correspondence, a perfect or analytically controlled target,
a determinant-line comparison, conductor and archimedean normalization,
and positivity for that same arithmetic function. A derived equivalence
or an AR translation does not supply a Hilbert-space adjoint, self-adjoint
domain or trace-class convergence. These remain the spectral program's
independent load-bearing obligations.

The executable calculation is \path{examples/coverings/theta_cover.py}.
The exact matrix obstruction is checked in \path{lean-analytic/CoveringBand.lean}.
The latter verifies selected actual matrices, not every category or
equivalence discussed here.

\section{From a Hom decomposition to a semi-invariant calculation}
For a finite-dimensional quotient, a specified projective presentation
$P_1\xrightarrow{d}P_0$ can define a determinant semi-invariant by applying
$\operatorname{Hom}(-,V)$, provided the resulting vector spaces have equal
dimension and bases or determinant lines have been fixed. The covering Hom
decomposition makes this matrix computable by deck components; it does not
make it block diagonal. Cross-degree terms matter. Already
\[
 \det\begin{pmatrix}1&a\\b&1\end{pmatrix}=1-ab
\]
differs from the product of its two diagonal entries. A product over deck
blocks is justified only after proving that the relevant differential or
endomorphism preserves those blocks. AR mesh relations concern morphism
compositions; a presentation of a ring of semi-invariants requires additional
generation and relation proofs, as in the earlier determinant chapters.
Neither the GLS root lattice nor a surface intersection form provides the
missing positive Hilbert inner product or arithmetic factor identity by
itself.
